% ============================================================================
%  ESTADÍSTICA INFERENCIAL — Sesión 5
%  Semana 5: 7 al 12 de septiembre de 2026
%  1.7 Muestreo probabilístico · 1.8 Muestreo no probabilístico
%  1.9 Cálculo del tamaño de la muestra · 1.10 Aplicaciones
%  Cierre de la Unidad 1 y repaso para el Examen 1
% ============================================================================
\documentclass[aspectratio=169,10pt]{beamer}

\usepackage[utf8]{inputenc}
\IfFileExists{spanish.ldf}{\usepackage[spanish,es-noshorthands]{babel}}{}
\usepackage{amsmath,amssymb}
\usepackage{booktabs}
\usepackage{tikz}
\usepackage{listings}

\usetheme{Uniminuto}

\uniasignatura{Estadística inferencial}
\uniprofesor{Andrés Felipe Puerta Velez}
\unicorreo{andres.puerta-v@uniminuto.edu.co}
\uninrc{CBASBL061}
\unisemana{Sesión 5 — Semana 5 · 7 al 12 de septiembre de 2026}

\newcommand{\listatiny}{\setbeamerfont{itemize/enumerate body}{size=\tiny}\tiny}

\newcommand{\fx}[1]{{\ttfamily\color{uniAzul}#1}}

% ------------------------------------------------------------- CÓDIGO R ----
\lstdefinestyle{Rcode}{
  language=R,
  basicstyle=\ttfamily\scriptsize,
  keywordstyle=\color{uniAzul}\bfseries,
  commentstyle=\color{uniGris}\itshape,
  stringstyle=\color{uniTitulo},
  showstringspaces=false, numbers=none, breaklines=true,
  columns=fullflexible,
  backgroundcolor=\color{uniCremaClaro},
  frame=leftline, framerule=1.2pt, rulecolor=\color{uniNaranja},
  xleftmargin=0.35cm, framexleftmargin=0.35cm,
  framexrightmargin=0.15cm, framextopmargin=2pt, framexbottommargin=2pt,
  aboveskip=0.15cm, belowskip=0.15cm
}
\lstdefinestyle{Rmini}{
  language=R,
  basicstyle=\ttfamily\tiny,
  keywordstyle=\color{uniAzul}\bfseries,
  commentstyle=\color{uniGris}\itshape,
  stringstyle=\color{uniTitulo},
  showstringspaces=false, numbers=none, breaklines=true,
  columns=fullflexible,
  backgroundcolor=\color{uniCremaClaro},
  frame=leftline, framerule=1.2pt, rulecolor=\color{uniNaranja},
  xleftmargin=0.35cm, framexleftmargin=0.35cm,
  framexrightmargin=0.15cm, framextopmargin=2pt, framexbottommargin=2pt,
  aboveskip=0.12cm, belowskip=0.12cm
}
\lstdefinestyle{Rout}{
  basicstyle=\ttfamily\tiny\color{black},
  showstringspaces=false, numbers=none, breaklines=true,
  columns=fullflexible,
  backgroundcolor=\color{uniGris!10},
  frame=leftline, framerule=1.2pt, rulecolor=\color{uniGris},
  xleftmargin=0.35cm, framexleftmargin=0.35cm,
  framexrightmargin=0.15cm, framextopmargin=2pt, framexbottommargin=2pt,
  aboveskip=0.12cm, belowskip=0.12cm
}

\begin{document}

\portada

\begin{frame}{Orden del día}
  \begin{enumerate}
    \item Llamado a lista y revisión de los ejercicios propuestos
    \item Repaso: dónde quedamos
    \item \azul{Temas 1.7 y 1.8: muestreo probabilístico y no probabilístico}
    \item \azul{Tema 1.9: cálculo del tamaño de la muestra}
    \item \azul{Tema 1.10: aplicaciones del muestreo en el entorno}
    \item Ejercicios en clase y \destacar{repaso para el Examen 1}
  \end{enumerate}
\end{frame}

\begin{frame}{Repaso: dónde quedamos}
  En la sesión 4 distinguimos \azul{parámetro} (población, fijo) de
  \azul{estadístico} (muestra, aleatorio), y vimos que una muestra grande mal
  elegida es peor que una pequeña bien elegida.

  \vspace{0.25cm}
  \begin{alertblock}{Las dos preguntas de hoy}
    \footnotesize
    \begin{enumerate}
      \item \destacar{¿Cómo} se elige la muestra? $\rightarrow$ técnicas de
            muestreo.
      \item \destacar{¿Cuántos} elementos se necesitan? $\rightarrow$ cálculo
            del tamaño de la muestra.
    \end{enumerate}
  \end{alertblock}

  \vspace{0.15cm}
  {\footnotesize\color{uniGris} Con esto se cierra la unidad 1. El
  \destacar{Examen 1} es la próxima semana: 14 al 19 de septiembre.}
\end{frame}

\seccion{1.7 y 1.8 Técnicas de muestreo}

\begin{frame}{Muestreo probabilístico}
  Todo elemento de la población tiene una \destacar{probabilidad conocida y
  mayor que cero} de ser seleccionado. Solo así se puede calcular el error
  muestral y hacer inferencia.

  \vspace{0.15cm}
  \scriptsize
  \begin{center}
  \begin{tabular}{@{}p{2.5cm}p{4.5cm}p{4.5cm}@{}}
    \toprule
    \textbf{Técnica} & \textbf{En qué consiste} & \textbf{Cuándo conviene} \\
    \midrule
    Aleatorio simple & Cada elemento con igual probabilidad; sorteo o números
      aleatorios & Población pequeña y homogénea, con marco muestral completo \\
    \addlinespace
    Sistemático & Se elige un arranque al azar y luego cada $k$-ésimo elemento,
      con $k=N/n$ & Listas ordenadas, líneas de producción \\
    \addlinespace
    Estratificado & Se divide en estratos homogéneos y se muestrea dentro de
      cada uno & La población tiene grupos claramente distintos \\
    \addlinespace
    Por conglomerados & Se divide en grupos, se sortean grupos completos y se
      observan enteros & Población dispersa geográficamente; abarata el trabajo \\
    \bottomrule
  \end{tabular}
  \end{center}
\end{frame}

\begin{frame}{Estratificado frente a conglomerados}
  \begin{columns}[T,onlytextwidth]
    \begin{column}{0.48\textwidth}
      \begin{block}{Estratificado}
        \footnotesize
        Los estratos son \azul{homogéneos por dentro} y distintos entre sí.
        Se toma una parte de \emph{cada} estrato.

        \vspace{0.15cm}
        \emph{Ejemplo:} clientes por ciudad; se encuesta en todas las ciudades.
      \end{block}
    \end{column}
    \begin{column}{0.48\textwidth}
      \begin{exampleblock}{Por conglomerados}
        \footnotesize
        Los conglomerados son \azul{heterogéneos por dentro} y parecidos entre
        sí. Se eligen \emph{algunos} y se observan completos.

        \vspace{0.15cm}
        \emph{Ejemplo:} se sortean 8 sucursales y se encuesta a todos sus
        clientes.
      \end{exampleblock}
    \end{column}
  \end{columns}

  \vspace{0.25cm}
  \begin{alertblock}{Afijación proporcional}
    \footnotesize
    En el estratificado, lo habitual es repartir $n$ en proporción al tamaño de
    cada estrato: $n_i = n\cdot N_i/N$. Con $N=1000$ repartido en 600, 300 y 100
    y $n=200$, corresponden \destacar{120, 60 y 20}.
  \end{alertblock}
\end{frame}

\begin{frame}{Muestreo no probabilístico}
  No se conoce la probabilidad de selección. \destacar{No permite calcular
  error muestral} ni generalizar con rigor, pero es útil en estudios
  exploratorios.

  \vspace{0.2cm}
  \footnotesize
  \begin{center}
  \begin{tabular}{@{}p{2.8cm}p{8.4cm}@{}}
    \toprule
    \textbf{Técnica} & \textbf{En qué consiste} \\
    \midrule
    Por conveniencia & Se toma lo que está a la mano: los que pasan por la
      entrada, los que contestan el formulario. \\
    Por cuotas & Se fijan cupos por grupo (edad, sexo) y el encuestador los
      llena como pueda. \\
    Por juicio o intencional & El investigador elige los casos que considera
      representativos o informativos. \\
    Bola de nieve & Cada participante refiere a otros. Útil en poblaciones
      difíciles de ubicar. \\
    \bottomrule
  \end{tabular}
  \end{center}

  \vspace{0.15cm}
  {\scriptsize\color{uniGris} Si un estudio reporta ``margen de error'' pero la
  muestra fue por conveniencia, ese margen no significa nada\,[2].}
\end{frame}

\seccion{1.9 Cálculo del tamaño de la muestra}

\begin{frame}{Fórmulas del tamaño de muestra para una proporción}
  \begin{columns}[T,onlytextwidth]
    \begin{column}{0.48\textwidth}
      \begin{block}{Población infinita o muy grande}
        \[ n=\frac{z^{2}\,p\,q}{e^{2}},\qquad q=1-p \]
      \end{block}
    \end{column}
    \begin{column}{0.48\textwidth}
      \begin{block}{Población finita de tamaño $N$}
        \[ n=\frac{N\,z^{2}\,p\,q}{e^{2}(N-1)+z^{2}\,p\,q} \]
      \end{block}
    \end{column}
  \end{columns}

  \vspace{0.25cm}
  \footnotesize
  \begin{center}
  \begin{tabular}{@{}llll@{}}
    \toprule
    \textbf{Símbolo} & \textbf{Qué es} & \textbf{Valor típico} & \textbf{Nota} \\
    \midrule
    $z$ & Valor crítico normal & 1{,}645 / 1{,}96 / 2{,}576 & para 90\% /
      95\% / 99\% \\
    $p$ & Proporción esperada & 0{,}5 si no se conoce & $p=0{,}5$ maximiza $n$ \\
    $e$ & Error máximo admitido & 0{,}03 a 0{,}05 & en puntos de proporción \\
    \bottomrule
  \end{tabular}
  \end{center}

  \vspace{0.15cm}
  {\scriptsize\color{uniGris} El resultado \emph{siempre} se redondea hacia
  arriba: 384{,}16 significa 385 encuestas.}
\end{frame}

\begin{frame}{Cómo se comporta $n$}
  \begin{columns}[T,onlytextwidth]
    \begin{column}{0.48\textwidth}
      \begin{exampleblock}{Al 95\% con $p=0{,}5$}
        \footnotesize
        \begin{center}
        \begin{tabular}{@{}cr@{}}
          \toprule
          Error $e$ & $n$ \\
          \midrule
          5\% & 385 \\
          3\% & 1\,068 \\
          2\% & 2\,401 \\
          \bottomrule
        \end{tabular}
        \end{center}
      \end{exampleblock}
    \end{column}
    \begin{column}{0.48\textwidth}
      \begin{alertblock}{Lecciones}
        \footnotesize
        \begin{itemize}
          \item Reducir el error a la mitad \azul{cuadruplica} la muestra:
                $n$ depende de $1/e^{2}$.
          \item Más confianza (99\%) exige más muestra.
          \item El tamaño de la población casi no importa cuando $N$ es grande:
                de $N=2\,500$ a $N=\infty$, $n$ pasa de 334 a 385.
        \end{itemize}
      \end{alertblock}
    \end{column}
  \end{columns}
\end{frame}

\begin{frame}{Ejercicio 1 — en clase}
  \begin{alertblock}{Individual (8 minutos)}
    Una cadena quiere estimar la proporción de clientes satisfechos, con
    \destacar{95\% de confianza} ($z=1{,}96$) y un error máximo del
    \destacar{5\%}. No hay estudios previos, así que se toma $p=0{,}5$.

    \vspace{0.25cm}
    \begin{enumerate}
      \item Calcule $n$ suponiendo población muy grande.
      \item Recalcule si la población es de $N=1\,200$ clientes.
      \item ¿Cuánto cambia $n$ si el error exigido baja al 3\%?
      \item Con $N=1\,200$ y $n=292$, ¿cada cuántos elementos se tomaría uno en
            un muestreo sistemático?
    \end{enumerate}
  \end{alertblock}
\end{frame}

\begin{frame}{Ejercicio 1 — solución}
  \small
  \begin{block}{Cálculos}
    \begin{align*}
      \text{1.}\;& n=\frac{1{,}96^{2}(0{,}5)(0{,}5)}{0{,}05^{2}}
                    =\frac{0{,}9604}{0{,}0025}=384{,}16 \Rightarrow \azul{385}\\[0.3em]
      \text{2.}\;& n=\frac{1200(0{,}9604)}{0{,}0025(1199)+0{,}9604}
                    =\frac{1152{,}48}{3{,}9579}=291{,}18 \Rightarrow \azul{292}\\[0.3em]
      \text{3.}\;& n=\frac{0{,}9604}{0{,}0009}=1067{,}11 \Rightarrow \azul{1068}
    \end{align*}
  \end{block}

  \begin{exampleblock}{Punto 4}
    \footnotesize
    $k=N/n=1200/292\approx 4$: arranque al azar entre 1 y 4 y luego uno de cada 4.
  \end{exampleblock}
\end{frame}

\begin{frame}{Ejercicio 2 — en clase}
  \begin{alertblock}{En parejas (8 minutos)}
    Una empresa tiene \destacar{4\,500 empleados} distribuidos así: 2\,700 en
    Operaciones, 1\,350 en Administración y 450 en Dirección. Se hará una
    encuesta de clima con \destacar{$n=300$}.

    \vspace{0.25cm}
    \begin{enumerate}
      \item ¿Cuántos empleados se encuestarían en cada área con afijación
            proporcional?
      \item ¿Qué técnica de muestreo se está usando y por qué es preferible al
            aleatorio simple en este caso?
      \item Si en lugar de eso se encuestara a los primeros 300 que respondan
            un correo, ¿qué tipo de muestreo sería y qué problema tendría?
    \end{enumerate}
  \end{alertblock}
\end{frame}

\begin{frame}{Ejercicio 2 — solución}
  \begin{columns}[T,onlytextwidth]
    \begin{column}{0.46\textwidth}
      \footnotesize
      \begin{center}
      \begin{tabular}{@{}lrrr@{}}
        \toprule
        Área & $N_i$ & Peso & $n_i$ \\
        \midrule
        Operaciones    & 2\,700 & 0{,}60 & \azul{180} \\
        Administración & 1\,350 & 0{,}30 & \azul{90} \\
        Dirección      &   450  & 0{,}10 & \azul{30} \\
        \midrule
        Total          & 4\,500 & 1{,}00 & 300 \\
        \bottomrule
      \end{tabular}
      \end{center}
    \end{column}
    \begin{column}{0.50\textwidth}
      \begin{block}{Punto 2}
        \footnotesize
        Es \azul{muestreo estratificado} con afijación proporcional. Garantiza
        que Dirección quede representada: con aleatorio simple podría no salir
        ningún directivo, porque son solo el 10\%.
      \end{block}
      \begin{alertblock}{Punto 3}
        \footnotesize
        Sería \destacar{por conveniencia}. Responden los más motivados o los
        más inconformes: el resultado queda sesgado y no se puede calcular
        error muestral.
      \end{alertblock}
    \end{column}
  \end{columns}
\end{frame}

\seccion{Repaso para el Examen 1}

\begin{frame}{Mapa de la Unidad 1}
  \scriptsize
  \begin{center}
  \begin{tabular}{@{}llp{6.6cm}@{}}
    \toprule
    \textbf{Tema} & \textbf{Sesión} & \textbf{Lo esencial} \\
    \midrule
    Conjuntos y eventos & 1 & $\cup$, $\cap$, $A^{c}$, $A\setminus B$,
      $A\triangle B$; excluyentes; $\#(A\cup B)=\#A+\#B-\#(A\cap B)$ \\
    Conteo & 2 & Multiplicación; $n!$; $P(n,r)$; $n^{r}$; $C(n,r)$;
      permutaciones con repetidos \\
    Probabilidad & 3 & Laplace, enfoques, axiomas, complemento, regla de la suma \\
    Condicional & 4 & $P(A\mid B)$, regla de la multiplicación, independencia,
      Bayes \\
    Muestreo & 4--5 & Parámetro vs. estadístico; técnicas; tamaño de muestra \\
    \bottomrule
  \end{tabular}
  \end{center}

  \vspace{0.2cm}
  \begin{alertblock}{Examen 1 — semana 6}
    \footnotesize
    \destacar{14 al 19 de septiembre}, en Aulas Virtuales (Moodle), presencial
    en el aula, individual, con tiempo límite. Vale el \destacar{33,3\%}.
    Se permite calculadora y formulario.
  \end{alertblock}
\end{frame}

\begin{frame}{Errores frecuentes que conviene evitar}
  \small
  \begin{itemize}
    \item Usar Laplace cuando los resultados \azul{no} son equiprobables (la
          suma de dos dados).
    \item Sumar probabilidades sin descontar la intersección.
    \item Confundir \azul{excluyentes} con \azul{independientes}.
    \item Usar $P(n,r)$ cuando el orden no importa, o $C(n,r)$ cuando sí.
    \item Olvidar que sin reemplazo las extracciones son dependientes.
    \item Redondear $n$ hacia abajo, o usar la fórmula infinita cuando $N$ es
          pequeño.
    \item Invertir el condicionamiento: $P(A\mid B)$ no es $P(B\mid A)$.
  \end{itemize}
\end{frame}

\begin{frame}[fragile]{Muestreo y tamaño de muestra, en R}
  \scriptsize
\begin{lstlisting}[style=Rmini]
set.seed(2026)
marco <- 1:1200                       # marco muestral de N = 1200

sample(marco, 60)                     # aleatorio simple
seq(sample(1:20, 1), 1200, by = 20)   # sistematico, k = N/n = 20

# Estratificado: 60 unidades repartidas proporcionalmente
estrato <- rep(c("A", "B", "C"), times = c(600, 400, 200))
tapply(marco, estrato, function(u) sample(u, round(length(u) * 60 / 1200)))

# Tamano de muestra para proporcion, poblacion infinita
n0 <- function(p, e, conf = 0.95) ceiling(qnorm(1 - (1 - conf)/2)^2 * p * (1 - p) / e^2)
n0(0.5, 0.05)                          # el caso mas conservador
ceiling(n0(0.5, 0.05) / (1 + (n0(0.5, 0.05) - 1) / 1200))   # con correccion
\end{lstlisting}
\begin{lstlisting}[style=Rout]
[1] 385
[1] 292
\end{lstlisting}

  \vspace{-0.1cm}
  {\tiny\color{uniGris} Escribir la fórmula como \destacar{función} permite
  probar en segundos qué pasa si el error admisible baja al 3\% o si la
  confianza sube al 99\%: se cambia un argumento, no toda la cuenta.}
\end{frame}

\begin{frame}{Cierre de la sesión}
  \begin{block}{Lo que debe quedar claro hoy}
    \begin{itemize}
      \item Solo el muestreo probabilístico permite calcular error e inferir.
      \item Estratificado: parte de cada grupo. Conglomerados: algunos grupos
            completos.
      \item $n=z^{2}pq/e^{2}$, con corrección por población finita si $N$ es
            pequeño; siempre se redondea hacia arriba.
      \item $n$ depende de $1/e^{2}$: exigir el doble de precisión cuesta
            cuatro veces más muestra.
    \end{itemize}
  \end{block}

  \vspace{0.2cm}
  \begin{alertblock}{Para el Examen 1 (semana 6)}
    \footnotesize
    Resolver los \destacar{20 ejercicios de la lámina siguiente} y repasar los
    de las sesiones 1 a 4: son el mejor simulacro disponible.
  \end{alertblock}
\end{frame}

\begin{frame}{Ejercicios propuestos}
  \vspace{-0.1cm}
  \begin{beamercolorbox}[wd=\textwidth,sep=0.10cm]{block body}
    \centering\scriptsize
    $z=1{,}645$ (90\%), $1{,}96$ (95\%), $2{,}576$ (99\%). Redondee hacia arriba.
  \end{beamercolorbox}
  \listatiny
  \vspace{0.15cm}
  \begin{columns}[T,onlytextwidth]
    \begin{column}{0.47\textwidth}
      \begin{enumerate}\setlength{\itemsep}{0.16cm}
        \item Infinita, 95\%, $p=0{,}5$, $e=0{,}05$
        \item Infinita, 95\%, $p=0{,}5$, $e=0{,}03$
        \item Infinita, 95\%, $p=0{,}5$, $e=0{,}02$
        \item Infinita, 99\%, $p=0{,}5$, $e=0{,}05$
        \item Infinita, 90\%, $p=0{,}5$, $e=0{,}05$
        \item Infinita, 95\%, $p=0{,}30$, $e=0{,}05$
        \item Infinita, 95\%, $p=0{,}80$, $e=0{,}05$
        \item Finita, $N=800$, 95\%, $e=0{,}05$
        \item Finita, $N=2\,500$, 95\%, $e=0{,}05$
        \item Finita, $N=150$, 95\%, $e=0{,}05$
      \end{enumerate}
    \end{column}
    \begin{column}{0.50\textwidth}
      \begin{enumerate}\setlength{\itemsep}{0.16cm}
        \setcounter{enumi}{10}
        \item Sistemático, $N=2\,000$, $n=200$: valor de $k$
        \item Sistemático, $N=1\,500$, $n=100$: valor de $k$
        \item Estratos 600/300/100, $n=200$: reparto
        \item Estratos 2\,700/1\,350/450, $n=300$: reparto
        \item Técnica: se sortean 8 sucursales completas
        \item Técnica: se encuesta a quien pase por la entrada
        \item Técnica: cada participante refiere al siguiente
        \item ¿Por qué se usa $p=0{,}5$ sin información previa?
        \item Si $e$ pasa de 0{,}05 a 0{,}025, ¿cuánto crece $n$?
        \item ¿Hay margen de error con muestreo por conveniencia?
      \end{enumerate}
    \end{column}
  \end{columns}
\end{frame}

\begin{frame}{Respuestas}{Autoevaluación de los ejercicios propuestos}
  \listatiny
  \vspace{0.15cm}
  \begin{columns}[T,onlytextwidth]
    \begin{column}{0.47\textwidth}
      \begin{enumerate}\setlength{\itemsep}{0.16cm}
        \item $384{,}16\Rightarrow 385$
        \item $1067{,}11\Rightarrow 1\,068$
        \item $2401\Rightarrow 2\,401$
        \item $663{,}58\Rightarrow 664$
        \item $270{,}60\Rightarrow 271$
        \item $322{,}69\Rightarrow 323$
        \item $245{,}86\Rightarrow 246$
        \item $259{,}86\Rightarrow 260$
        \item $333{,}02\Rightarrow 334$
        \item $108{,}09\Rightarrow 109$
      \end{enumerate}
    \end{column}
    \begin{column}{0.50\textwidth}
      \begin{enumerate}\setlength{\itemsep}{0.16cm}
        \setcounter{enumi}{10}
        \item $k=10$
        \item $k=15$
        \item $120$, $60$ y $20$
        \item $180$, $90$ y $30$
        \item Por conglomerados
        \item Por conveniencia
        \item Bola de nieve
        \item Porque $pq$ es máximo en $p=0{,}5$: da el $n$ más grande y por
              tanto seguro
        \item Por $4$ \;{\color{uniTitulo}($n\propto 1/e^{2}$)}
        \item No: la probabilidad de selección es desconocida
      \end{enumerate}
    \end{column}
  \end{columns}
\end{frame}

\begin{frame}{Referencias}
  \footnotesize
  \begin{itemize}
    \item[{[1]}] Díaz Rodríguez, M. (2022). \textit{Estadística inferencial
      aplicada}. Universidad del Norte.
      \texttt{https://elibro.net/es/lc/uniminuto/titulos/221614}
    \item[{[2]}] Llinás Solano, H. (2014). \textit{Introducción a la
      estadística con aplicaciones en ciencias sociales}. Universidad del
      Norte. \texttt{https://elibro.net/es/lc/uniminuto/titulos/70062}
    \item[{[3]}] Contento Rubio, M. R. (2019). \textit{Estadística con
      aplicaciones en R}. Editorial Utadeo.
      \texttt{https://elibro.net/es/lc/uniminuto/titulos/220926}
  \end{itemize}

  \vspace{0.3cm}
  \begin{center}
    {\Large\color{uniAzul}\textbf{Éxitos en el Examen 1}}
  \end{center}
\end{frame}

\end{document}
